\documentclass[10pt]{article}
\usepackage[margin=2.2cm]{geometry}
\usepackage[T1]{fontenc}
\usepackage{lmodern}
\usepackage{textcomp}
\usepackage{amsmath,amssymb}
\usepackage{enumitem}
\usepackage{xcolor}
\usepackage[hidelinks]{hyperref}
\setlength{\parindent}{0pt}
\setlength{\parskip}{4pt}
\newcommand{\rev}[1]{\par\medskip\noindent\textbf{Reviewer comment #1}\par}
\newcommand{\resp}{\par\noindent\textcolor{blue!50!black}{\textbf{Response.}}~}
\newcommand{\chg}{\par\noindent\textcolor{green!35!black}{\textbf{Changes.}}~}

\begin{document}

\begin{center}
{\Large\bfseries Response to the Reviewer}\\[4pt]
Manuscript: \emph{Corner-Speed-Based Design and Optimization of a Switched Reluctance Motor Drive for Low-Voltage Electric Bicycles}\\
(previous title: \emph{Analysis and Design Considerations for a Switched Reluctance Motor-Based Drive System for an Electric Bike})
\end{center}

We thank the reviewer for a detailed and constructive review. We agree with
the main criticisms, and the revision substantially extends the work. The
largest changes are as follows.
\begin{enumerate}[nosep]
\item \textbf{FE model.} A nonlinear 2-D finite-element (FE) model of the full
machine was added (Section~IV-B). It includes a mesh-convergence study and a
quantitative comparison with the analytical model at 8, 15, 23 and 31~A
(Table~VI, Fig.~3).
\item \textbf{FE-based results.} All drive-level results are now computed with
the FE flux-linkage characteristic instead of the analytical one.
\item \textbf{TSF control.} Torque-sharing-function control (linear, cubic and
exponential) was implemented and compared with current-reference control
(Section~VII).
\item \textbf{Sensitivity studies.} New studies cover PWM frequency, a
four-node transient thermal network, the gear model, DC-link current sharing,
air gap, pole arcs, the dwell limit and the convergence of the angle search.
\item \textbf{Novelty.} A new subsection (II-E) positions $P_c=k_cV_{dc}I_{pk}$
relative to the classical volt-second and converter volt-ampere arguments.
\end{enumerate}
\textbf{Experimental validation has not been added.} We state this explicitly
in the abstract, the introduction and Section~X-D, and we reframed the paper as
an analytical, FE and simulation-based design methodology (see comment~5).

Because the FE characteristic and a more thorough angle search replaced the
analytical model, several headline numbers changed. Table~1 below lists them.
Line-by-line changes are given under each comment.

\begin{center}
\small
\begin{tabular}{lcc}
\hline
Quantity & Original manuscript & Revised (FE-based)\\
\hline
Current for launch torque / converter limit & 27.9 A / 30.7 A & 26.5 A / 29.2 A\\
Launch torque at limit (margin) & 6.81 N$\cdot$m (+14.9\%) & 6.63 N$\cdot$m (+11.6\%)\\
Rated $\eta_{mot}$ / $\eta_{sys}$ & 75.8 / 72.8\% & 80.2 / 77.9\%\\
Peak $\eta_{mot}$ / $\eta_{sys}$ (map) & 81.4 / 78.6\% & 80.6 / 78.3\%\\
Torque ripple, launch / rated (current-ref.) & 113 / 148\% & 60 / 87\%\\
Route consumption / estimated range & 5.05 Wh/km / 64 km & 5.12 Wh/km / 63 km\\
\hline
\multicolumn{3}{l}{\footnotesize Table 1: Main numerical changes. The rated-efficiency increase comes mostly from the improved angle search (comment 26).}
\end{tabular}
\end{center}

%-----------------------------------------------------------------------------
\rev{2--3 (strengths; novelty of the corner-speed argument)}
The reviewer asks us to establish whether $P_c=k_cV_{dc}I_{pk}$ is new or a
reformulation of established volt-second and energy relationships.
\resp We agree that the relation is built from established ingredients. The
volt-second limit is standard. The energy-conversion-loop argument that links
converter volt-amperes to output power through the $\psi$--$i$ loop
utilization is classical (Miller's energy ratio and converter volt-ampere
requirement). We no longer present the relation as a new law. The
contribution is now stated as how the relation is \emph{used} for low-voltage
sizing:
\begin{itemize}[nosep]
\item choosing the turns at a mission-derived corner speed rather than at the
speed limit;
\item showing that the converter current is fixed before the machine is sized,
so that enlarging the machine buys efficiency, not current;
\item reporting $k_c$ for this machine class.
\end{itemize}
\chg New Section~II-E, \emph{Relationship to Existing Volt-Second and
Energy-Based Methods}, as the reviewer suggested. Contribution~1 in the
introduction and the conclusion are reworded accordingly. The trade study
confirms the claim quantitatively: $I_{pk}$ varies by 22\% over a fourfold
range of mass, with $k_c=0.54$--$0.66$ across nine designs.

\rev{4 (electromagnetic model insufficiently validated; FE required)}
\resp Done. We implemented a nonlinear 2-D magnetostatic FE model of the full
12/8 cross-section with the following features:
\begin{itemize}[nosep]
\item vector potential formulation, first-order triangles and Newton iteration;
\item M270-35A B--H curve and re-meshing at every rotor position;
\item Arkkio torque, cross-checked against the co-energy derivative to within
0.5\%.
\end{itemize}
Mesh convergence is shown in Table~V: halving the gap element size changes
results by less than 0.1\%. As requested, Table~VI and Fig.~3 compare the
analytical model and FE at 8, 15, 23 and 31~A:
\begin{itemize}[nosep]
\item flux linkage: RMS error 13--15\%, maximum 23\%;
\item aligned inductance: $+5.7$\%; unaligned inductance: $-9.7$\%;
\item mean torque: $-2$ to $-17$\%; peak torque: within $\pm$10\%.
\end{itemize}
Because these errors are significant, \emph{all drive-level results now use the
FE characteristic}. It is imported as $\psi(i,\theta)$ tables, with torque
computed from the co-energy of the same table so that the model stays
energy-consistent. Fig.~4 shows the FE field solution.
\chg Sections~IV-B and IV-C, Tables~V and VI, Figs.~3 and 4. Sections~VI--VIII
were recomputed with the FE characteristic.

\rev{5 (no experimental prototype)}
\resp We fully agree that experimental validation would strengthen the paper.
It is not available for this revision, and we have not presented any
simulated quantity as measured. Following the reviewer's alternative
recommendation, the paper has been reframed.
\begin{itemize}[nosep]
\item The title now emphasizes the design methodology.
\item The abstract states that ``No prototype measurements are reported, and
all performance figures are model predictions''.
\item The introduction limits the scope explicitly.
\item Section~X-D lists the planned dynamometer tests (phase voltage and
current, torque and ripple, DC input power, efficiency map, temperature rise,
acoustic noise).
\end{itemize}
\chg Title, abstract, Section~I (scope paragraph), Section~X-D and the conclusion.
\emph{[Note to authors: if prototype data become available before
resubmission, Section~X-D is where they would replace the stated limitation.]}

\rev{6 (torque ripple; TSF only discussed, not implemented)}
\resp Done. We implemented linear, cubic and exponential TSFs with current
references obtained from the inverse FE torque characteristic and tracked by a
three-level hysteresis controller. Table~XII and Fig.~11 compare them with
current-reference control (CCC), in the format the reviewer proposed:
\begin{itemize}[nosep]
\item At 200~r/min, cubic TSF reduces ripple from 197\% to 10\%.
\item At 500~r/min it reaches 33\%, and at 800~r/min and 1.5~N$\cdot$m, 39\%
(against 83\% for CCC).
\item The cost is 12--28\% higher RMS current and 7--12 points of efficiency.
\item Above about 800~r/min, no TSF solution exists within $I_{max}$ at 36~V.
\end{itemize}
This is the known speed limit of TSF, and it is particularly restrictive at
low voltage. We therefore recommend a hybrid controller: TSF for launch,
walk-assist and low-speed climbing, and angle control above about 12~km/h. The
FE-based current-reference ripple is itself lower than originally predicted
(87\% at the rated point instead of 148\%).
\chg New Section~VII, Table~XII, Fig.~11, and the abstract and conclusion.

\rev{7 (launch margin should be stated)}
\resp Done. The FE-based capability at the current limit is 6.63~N$\cdot$m,
which exceeds the 5.93-N$\cdot$m requirement by 11.6\%. The original
analytical figure was 6.81~N$\cdot$m (+14.9\%). The gear-sensitivity study
shows that the margin shrinks to 6\% if $\eta_g=0.90$.
\chg Section~VI-A, abstract, Table~XIII.

\rev{8 (efficiency interpretation; comparison with PM motors)}
\resp We agree that the comparison was not like-for-like. We removed the
quantitative comparison with commercial PM hubs, and Section~X-C now states
why a fair comparison requires identical boundaries, voltage, temperature and
method. Efficiency boundaries are defined explicitly in Section~IV-F
($\eta_{mot}$, $\eta_{sys}$ and $\eta_{conv}$, with the gear excluded, so that
battery-to-wheel efficiency equals $\eta_{sys}\eta_g$). The Fig.~8 caption
states the boundary.
\chg Sections~IV-F and X-C, Fig.~8 caption.

\rev{9 (thermal model too simplified; sensitivity)}
\resp Done. The one-node model was replaced by a transient four-node network
(winding, stator, rotor, shell) with temperature-dependent copper loss.
Table~XI gives the sensitivity the reviewer requested: $h=15$/25/40~W/m$^2$K,
ambient 35/45~$^\circ$C (25~$^\circ$C is also computed), continuous 250~W,
6\% climb, and repeated launches (5~s every 60~s). Fig.~10 shows temperature
versus time. One new finding is that a sustained 6\% climb at
$h=15$~W/m$^2$K and 45~$^\circ$C reaches 142~$^\circ$C at steady state. We
therefore now recommend class-H insulation or thermal derating for hot or
mountainous use, instead of claiming a wide margin.
\chg Section~VI-F, Table~XI, Fig.~10.

\rev{10 (DC-link capacitor derivation)}
\resp Done. Section~VI-E now gives the charge bound explicitly as
Eq.~(16) and explains that it assumes the battery supplies only the mean
current. It gives 1.56~mF for 5\% ripple at the rated point, but a meaningless
30~mF at launch, where the stroke frequency is only 32~Hz. We therefore added
a harmonic current-sharing analysis between the capacitor
($R_{ESR}=20$~m$\Omega$) and the battery ($R_b=0.15~\Omega$, $L_b=1~\mu$H),
with all assumptions stated (Table~X). The conclusions are:
\begin{itemize}[nosep]
\item 1~mF gives 4.8\% ripple at the rated point;
\item the launch ripple is set by battery impedance;
\item the bank is governed by ripple current (13.4~A RMS), not capacitance.
\end{itemize}
The original 3.8-mF recommendation is withdrawn in favor of 1--2.2~mF rated
for at least 14~A RMS.
\chg Section~VI-E, Eq.~(16), Table~X.

\rev{11 (switching frequency; production PWM controller)}
\resp Done. We implemented fixed-frequency peak-current-mode PWM in the
simulator and accumulated switching energy per transition at the actual
current, including diode reverse recovery. The hill, rated and launch points
were evaluated at 4, 8, 16 and 20~kHz, as requested (Table~IX). At 36~V,
switching loss is at most 1.9~W (launch, 20~kHz), and $\eta_{sys}$ changes by
less than 0.1 point, so the efficiency results also hold for the recommended
production PWM.
\chg Section~VI-D, Table~IX.

\rev{12 (gearbox efficiency assumption)}
\resp Done. The gear model now includes an optional speed-proportional drag
loss in addition to the mesh efficiency (Eq.~(2)). Table~XIII reports the
sensitivity to $\eta_g=0.90$, 0.93, 0.95 and 0.97 and to two drag cases.
Consumption ranges from 5.02 to 5.38~Wh/km, and the launch requirement ranges
from 5.81 to 6.26~N$\cdot$m. The launch requirement stays within the envelope
in all cases.
\chg Section~II-B, Eq.~(2), Section~VIII, Table~XIII.

\rev{13 (route model quasi-static; braking assumption)}
\resp We now list the neglected effects explicitly: transients, battery sag,
controller delay and thermal state. Battery sag is partly covered by design,
because the turns are chosen at the 30-V end-of-discharge voltage. The
mechanical-braking assumption is now justified quantitatively. The brakes
dissipate 7.2~Wh on the route, so regeneration could return at most about
5~Wh (about 15\% of battery energy on this hilly route). The asymmetric bridge
supports regeneration without hardware change.
\chg Section~VIII.

\rev{14 (64-km range claim)}
\resp Changed to ``an \emph{estimated range of about 63~km under these route
and modeling assumptions}'' everywhere, with the gear sensitivity (60--65~km).
\chg Abstract, Section~VIII, Table~XIII, conclusion.

\rev{15 (position sensing details)}
\resp Expanded. The text now covers extrapolation with speed and acceleration,
quantization and sensor alignment ($\pm1^\circ$) as commutation-angle errors,
the stand-still hill start (exciting the phase pair that gives positive torque
throughout the 15$^\circ$ Hall sector), and direction reversal.
\chg Section~X-A.

\rev{16 (mechanical feasibility; tolerance stack-up)}
\resp Done in two parts. First, FE air-gap sensitivity (Table~XV): at 0.35~mm
the 30-A mean torque drops by 4.0\% (8.3\% at 15~A), and at 0.40~mm by 8.2\%
(16\%). Second, a tolerance allocation in the format the reviewer suggested
(Table~XVI), clearly labeled as an assumed design budget rather than a
measurement. It gives a minimum local gap of 0.215~mm worst case and 0.26~mm
by RSS, with the corresponding eccentricity and unbalanced magnetic pull.
\chg Section~X-B, Tables~XV and XVI.

\rev{17 (mutual coupling neglected)}
\resp Quantified by FE with two phases excited simultaneously (Section~IV-D).
In the commutation overlap, the second phase changes the flux linkage by at
most 7.2\%, and the torque differs from superposition by at most
0.36~N$\cdot$m (6.9\%). The mutual coupling is not yet included in the drive
model. We state this, and we note that it bounds the achievable TSF ripple
without online compensation.
\chg Section~IV-D, Sections~VII and X-D.

\rev{18 (literature review needs 2020--2026 references)}
\resp Added verified recent references:
\begin{itemize}[nosep]
\item Terzi\'c and Mihi\'c (2022), an SRM design for a mid-drive e-bike, the
closest prior work, which the paper now positions against;
\item Lan \emph{et al.} (2021), a review of SRM EV powertrains;
\item Fang \emph{et al.} (2021), a review of SRM control and vibration;
\item Abdalmagid \emph{et al.} (2022), geometry and topology optimization;
\item Hamouda \emph{et al.} (2020), torque control for light EVs;
\item Bostanci \emph{et al.} (IEEE TTE, 2017);
\item Xue \emph{et al.} (2009), TSF optimization.
\end{itemize}
\emph{[Note to authors: please add further 2023--2026 papers you consider
essential, e.g.\ on SRM acoustic noise and integrated hub motors; we only
included references whose bibliographic details we could verify.]}

\rev{19 (FE/experimental plots)}
\resp FE plots were added: the validation figure (Fig.~3) and the field plot
(Fig.~4). All performance figures (Figs.~6--9 and 12) are now FE-based. No
experimental plots are available; see comment~5.

\rev{20 (title)}
\resp Adopted the reviewer's preferred title: \emph{Corner-Speed-Based Design
and Optimization of a Switched Reluctance Motor Drive for Low-Voltage Electric
Bicycles}.

\rev{21 (summary of major comments)}
\resp Major comments 1--7 are addressed as described under comments~4--5
(validation), 6 (torque ripple), 2--3 (novelty), 9 and 4 (loss and thermal
models), 11 (controller mismatch), 12 (gear) and 14 (range).

\rev{23 (``halves the launch current'' phrasing)}
\resp Rephrased as the reviewer suggested. Section~II-D and the conclusion now
read: ``For the specified 36-V, 500-W corner power, this approximately halves
the peak current relative to sizing the turns for launch torque at 25~km/h''.
The assumptions (one-stroke dwell, fixed $V_{dc}$, similar $k_c$) are stated
next to the claim.

\rev{25 (minor corrections)}
\begin{itemize}[nosep]
\item \emph{Placeholder authors.} These must be filled in by the authors; they
are marked TODO in the source.
\item \emph{Mass definition.} 3.4~kg is now defined as the active
electromagnetic mass (laminations and copper), excluding housing, shaft,
bearings and gear (Section~III-D, Table~IV).
\item \emph{Efficiency definitions and Fig.~7 boundary.} See comment~8
(Section~IV-F, Fig.~8 caption).
\item \emph{1.1\% power-balance error.} With the FE tables it is 1.4\%. It is
now explained as the error of the explicit time integration and the bilinear
table interpolation (Section~IV-E).
\item \emph{Repository link.} The repository URL is given in the first-page
footnote and as reference [20].
\item \emph{Optimization grid and convergence.} Full details are in
Section~V-A, together with a convergence check against a 4$\times$ denser
grid: $\eta_{sys}$ improves by at most 1.5 points and maximum torque by at most
2.9\%.
\item \emph{Why $0.55\tau_r$.} A dwell sweep shows that $0.55\tau_r$ maximizes
base-speed torque (2.07, 2.76, 3.50, 3.95 and 3.77~N$\cdot$m for 0.40--0.60 and
above). The original text also wrongly stated that the limit avoids continuous
conduction. The optimum above base speed in fact runs in mild continuous
conduction (minimum current 1.3--1.7~A), and this is now stated (Section~V-A).
\item \emph{Pole-arc sensitivity.} Six FE pole-arc variants are reported in
Table~III.
\item \emph{Uncertainty.} Model-uncertainty statements are collected in
Section~X-D. Error bars are not meaningful without measurements, but the
FE-versus-analytical differences (Table~VI) and the grid-convergence bounds
quantify the model uncertainty.
\end{itemize}

\rev{26 (additional change made by the authors)}
\resp While addressing the dwell-limit question, we found that the original
angle search was too coarse. A tighter dwell limit produced a \emph{better}
rated point. We therefore added dwell-based turn-on candidates to the search
and verified convergence (comment~25). This, more than the change to the FE
characteristic, explains why the rated efficiency increased from 72.8\% to
77.9\%. The trade study was recomputed with the improved search. It no longer
supports the original claim that 11:1 gearing gains about four points at
equal mass: at similar mass, 8:1 and 11:1 are within one point. Section~IX has
been corrected accordingly.

\bigskip
All code, FE scripts and data needed to reproduce every table and figure are
in the public repository cited in the manuscript.

\end{document}
